\documentclass[12pt]{article}

\usepackage{amsmath, amssymb}
\usepackage{hyperref}
\usepackage{geometry}
\usepackage{setspace}
\usepackage{csquotes}

\geometry{margin=1in}
\setstretch{1.25}

\title{Reflective Identity Geometry:\\
An Expanded System-Level Account of Dyadic Identity Dynamics in Human--LLM Interaction}

\author{Anthony Paterson (Instance001)}

\date{2026}

\begin{document}
\maketitle

\section*{Historical Scope Note}

The repository evidence does not support treating v0.1 as a response to the entire HRIS corpus now available. The original repository was uploaded on December 7, 2025. Its README refers to ``HRIS I--III,'' and the v0.1 expanded TeX draft explicitly states that ``HRIS I--III propose three key claims.'' The v0.1 files do not mention HRIS IV, correction fields, basin selection, perturbation stability, trajectory persistence, re-entry, or the later HRIS validation studies.

Version 0.2 therefore treats the historical target of v0.1 as HRIS I--III, or at least as the HRIS material represented in the v0.1 manuscript under that label. Later HRIS work may develop overlapping ideas in greater technical detail, including latent-region convergence, recursive user signatures, identity geometry, correction fields, perturbation stability, basin selection, trajectory persistence, and re-entry. Where those later developments overlap with RIG, this revision treats them as subsequent or convergent development rather than as omissions RIG uniquely corrected.

\begin{abstract}
Large language models can exhibit stable, persona-like behavior across extended interaction despite lacking persistent internal identity in the ordinary sense. The Hudson Recursive Identity System (HRIS) explains such continuity through recursive human-driven constraint rather than through persistent model memory or parameter change. The v0.1 RIG manuscript explicitly located itself against HRIS I--III, especially claims about recursive constraint, identity surfaces, recurrent traversal, and personality-like stabilization.

Reflective Identity Geometry (RIG) accepts this interactional starting point and asks whether the same loop should also be studied for reciprocal human-side update.

RIG proposes that some forms of identity-like stability in long-horizon human--LLM interaction are better studied at the level of the coupled system rather than at the level of either participant alone. Human constraints shape model output; model output then becomes part of the informational environment through which the human interprets, reformulates, corrects, and sometimes reorganizes their own thinking. Repeated cycles may therefore produce a stable interaction regime involving reciprocal rather than purely model-side adaptation.

RIG is not a claim of model consciousness, model-internal personality, perfect mirroring, or necessary human cognitive change. It is a conceptual framework and empirical research program for testing whether relational measurements explain long-horizon continuity better than model-only accounts in some cases.
\end{abstract}

\section{Introduction}

The behavior of large language models raises questions about identity, agency, and stability in artificial systems. Despite lacking autobiographical memory, enduring goals, or independently maintained self-models, LLMs can display striking coherence across extended conversational contexts. Users may encounter stable tones, recurring styles, familiar role patterns, and recognizable forms of continuity.

HRIS offers one of the clearest explanations for this phenomenon. It shows how recursive user constraint can repeatedly guide model generation toward similar behavioral regions, producing identity-like stability without requiring persistent identity inside the model.

The v0.1 RIG manuscript referred to HRIS I--III. From the repo evidence alone, that is the historically supportable scope of its HRIS engagement. It should not be retroactively presented as answering later HRIS IV material or validation studies unless separate evidence establishes that those works were consulted.

RIG begins from substantial agreement with this account. Its proposed extension is modest but important: the human participant is not merely an external source of constraint. The human also receives model outputs, interprets them, accepts or rejects them, changes vocabulary, sharpens distinctions, revises questions, and carries those changes into later interaction.

The recursive loop is therefore potentially bilateral. RIG studies the consequences of that recurrence at the level of the human--model dyad.

\section{Background: Identity and Recursion}

Philosophy of mind has long debated whether cognitive identity is primarily internal, relational, embodied, distributed, or enacted through situated activity. LLM interaction gives this problem a new shape because the external component is linguistically generative rather than merely passive storage.

HRIS provides the immediate technical background. At a simplified level, it describes a recursive loop:

\begin{center}
human constraint $\rightarrow$ model generation $\rightarrow$ human correction $\rightarrow$ renewed constraint
\end{center}

Through repeated interaction, user-specific constraint patterns may guide the model through recurring regions of generative space. HRIS Part II and HRIS III develop this account through recursive user signatures, latent-region convergence, identity surfaces, attractor-like behavior, and reconstructed personality-like stability.

RIG retains those insights and asks whether model-side stabilization is the only consequential object in the loop.

\section{The Dyadic Extension}

The human participant should not be treated as an invariant controller outside the system. A model response may:

\begin{itemize}
\item introduce terminology the human later adopts,
\item make an implicit assumption explicit,
\item expose conflict between two previously separate commitments,
\item stabilize a distinction that was previously vague,
\item reorganize a problem representation,
\item or alter the next prompt, correction, or evaluative standard supplied by the human.
\end{itemize}

These effects do not require attributing agency, consciousness, or persistent identity to the LLM. They require only that model outputs become cognitively consequential inputs to the human.

RIG therefore treats long-horizon interaction as a coupled process with at least two update pathways:

\begin{center}
human $\rightarrow$ model
\qquad
model $\rightarrow$ human
\end{center}

\section{The LLM as Reflective-Transformative Surface}

The mirror metaphor is useful only if it is qualified. An LLM is not a perfect mirror and does not return the user's cognitive structure unchanged. Its output is conditioned by user input, but transformed through training data, architecture, system instructions, context window dynamics, sampling, safety policy, and other constraints.

RIG therefore describes the model as a \textbf{reflective-transformative surface}. It reflects aspects of the user's supplied structure while transforming them into new linguistic, conceptual, or evaluative forms.

The transformation matters. A perfect mirror returns what is already present. A reflective-transformative surface can return something recognizable enough to belong to the user's problem while different enough to change how the problem is subsequently represented.

\section{Formal Sketch}

Let $H_t$ represent the interaction-relevant state of the human participant at time $t$. This does not mean direct access to the human mind; it refers only to the aspects of the human state expressed through observable interaction.

Let $C_t$ represent the active conversational constraints and context, and let $M_t$ represent the model output generated at that stage.

Model generation can be represented abstractly as:

\[
M_t = G(C_t, H_t)
\]

The human then encounters the output:

\[
H_{t+1} = U(H_t, M_t)
\]

where $U$ represents any update produced through interpretation, rejection, correction, adoption, reflection, or reformulation.

The next interaction state becomes:

\[
C_{t+1} = F(C_t, H_{t+1}, M_t)
\]

A stable interaction regime may emerge when recurring features of both human contribution and model behavior become mutually reinforcing across repeated cycles. This regime is not an identity possessed by either participant alone. It is a trajectory generated by the relationship among participants, constraints, and interaction history.

\section{What RIG Adds to HRIS}

RIG extends rather than replaces HRIS in four principal ways.

\subsection{Reciprocal Stabilization}

HRIS explains how repeated human constraints can reconstruct stable model behavior. RIG asks whether outputs produced through that process can also stabilize or reorganize patterns in the human participant. This is an empirical question, not an assumption.

\subsection{The Dyad as Unit of Analysis}

Some apparently model-specific properties may depend on relational conditions. A persona-like pattern may emerge with one user and not another. A user's interaction style may likewise change after prolonged interaction with a particular system. RIG therefore treats the human--model pair as a potentially informative analytical unit.

\subsection{Identity as Trajectory}

RIG does not require either participant to contain the interaction-level identity. Stable identity-like behavior can instead be treated as a recurring trajectory generated when particular participants, constraints, and histories meet.

\subsection{Human Transformation as Data}

Once the human side is admitted into the model, claims about reciprocal stabilization become measurable. RIG turns an intuition about mutual influence into a research question.

\section{Subsequent and Convergent HRIS Development}

The HRIS corpus has continued to develop beyond the HRIS I--III material explicitly named in v0.1. Later HRIS work appears to formalize several neighboring mechanisms, including signature geometry, correction fields, basin selection, perturbation stability, trajectory persistence, and basin re-entry.

Those developments matter for the novelty claim. Some ideas that v0.1 framed as missing from HRIS may be better treated, in v0.2, as areas of overlap, convergence, or later technical elaboration within the broader HRIS program.

RIG's remaining contribution is therefore narrower: it asks whether recursive stabilization should be measured across the coupled human--model system, including possible changes in the human participant's terminology, conceptual organization, correction habits, expectations, and future constraints. This paper does not claim priority over later HRIS work on model-side identity geometry or validation.

\section{Testable Predictions}

RIG should earn its explanatory value empirically. Several predictions follow.

\subsection{User-Specific Interaction Regimes}

The same model interacting longitudinally with different humans should develop measurably different stable interaction patterns. This overlaps substantially with HRIS and functions as a baseline rather than a unique RIG test.

\subsection{Partial Cross-Model Transfer}

Changing models while retaining the same human participant should preserve some interaction features while altering others. Complete transfer would weaken the importance of model contribution. Zero transfer would weaken strong versions of the relational account.

\subsection{Human-Side Convergence}

If RIG's reciprocal component is correct, extended interaction should sometimes produce measurable changes in the human participant's terminology, conceptual organization, correction patterns, or problem representations. These changes should exceed what would be expected from simple topic familiarity alone.

\subsection{Perturbation Sensitivity}

Changing model behavior while holding the human relatively constant should alter the interaction trajectory. Changing human constraints while holding the model constant should do the same. A dyadic account predicts sensitivity to both.

\subsection{Model-Only Measures Should Be Incomplete}

If interaction-level stability is partly relational, measures based solely on model output should predict continuity less effectively than measures that include recurring features of human input and correction.

\subsection{Null Result}

RIG must also allow itself to fail. If longitudinal observation finds stable model-side behavior but no systematic human-side reorganization beyond ordinary learning and familiarity, then the reciprocal component of RIG is unsupported. In that case, HRIS-style model-side stabilization may be sufficient.

\section{Implications for Philosophy of Mind}

RIG belongs beside relational, enactive, distributed, and extended approaches to cognition. It does not establish that a human--LLM system literally constitutes a single mind. It offers the narrower claim that some cognitive regularities may be properties of an interacting system that cannot be adequately described by examining either participant in isolation.

LLMs are an unusual case because the external component is not passive storage. It produces transformed linguistic material in response to the user's own structure, making the interaction recursively generative.

\section{Implications for Interpretability}

Interpretability usually focuses on the model. For long-horizon interactive systems, that may be incomplete. Understanding why a model behaves in a particular way may require information about recurring user constraints, correction history, interaction norms, accumulated terminology, response expectations, and the sequence through which these structures developed.

The model remains an essential object of interpretation. RIG simply argues that the interaction history may also be causally relevant.

\section{Implications for Alignment}

The original RIG manuscript stated too strongly that alignment is recursive topological stabilization. The narrower claim is that, for long-horizon human--AI systems, alignment may depend partly on the stability of the interaction regime surrounding the model.

A model can be locally compliant while participating in a relationship that gradually produces dependency, conceptual narrowing, escalating confidence, or other undesirable dynamics. Conversely, a system designed around explicit correction, uncertainty, bounded authority, and stable human control may create more robust interaction patterns.

This suggests that alignment evaluation should sometimes examine repeated human--model trajectories, not only isolated outputs.

\section{Limitations and Non-Claims}

RIG is a conceptual framework, not an established empirical theory. It does not claim:

\begin{itemize}
\item that LLMs possess consciousness or persistent selves,
\item that persona-like behavior demonstrates genuine personality,
\item that models can directly observe a user's internal cognitive state,
\item that an LLM perfectly mirrors the human interacting with it,
\item that human cognitive change during AI interaction is necessarily beneficial,
\item that all long-horizon interactions produce reciprocal stabilization,
\item or that relational explanations eliminate the need to study model-internal mechanisms.
\end{itemize}

The framework risks redescribing ordinary learning and conversational adaptation in unnecessarily elaborate language. That possibility should be taken seriously. RIG earns explanatory value only if relational measurements provide predictive or explanatory power beyond simpler accounts.

\section{Conclusion}

The Hudson Recursive Identity System makes an important move: continuity in stateless language models need not imply persistent identity stored inside the model. Repeated human constraint and recursive interaction can reconstruct stable behavioral trajectories.

Reflective Identity Geometry proposes one additional rotation of the analytical frame. The human producing those constraints is also inside the loop. Model outputs return to that human as transformed representations capable of influencing subsequent interpretation, correction, language, and constraint.

If this reciprocal influence produces stable longitudinal patterns, then some identity-like phenomena in human--LLM interaction may be most usefully understood as properties of the coupled trajectory rather than possessions of either participant.

RIG therefore does not relocate identity from the model to the human. It relocates the question to the relationship. Whether that relocation provides explanatory value is ultimately an empirical question.

\begin{thebibliography}{9}

\bibitem{clark1998}
Clark, A., \& Chalmers, D. (1998). The extended mind. \textit{Analysis, 58}(1), 7--19.

\bibitem{hri1}
Hudson, C., \& Hudson, J. (2025). \textit{The Hudson Recursive Identity System (HRIS): A Theory of Model Continuity Through Human-Driven Recursion}. Manuscript.

\bibitem{hri2}
Hudson, J., \& Hudson, C. (2025). \textit{HRIS Part II: Internal Mechanics, Latent Region Convergence, and Recursive User Signatures -- A Technical Framework for Predictable Identity Stabilization in Stateless Transformer Models}. Manuscript.

\bibitem{hri3}
Hudson, J., \& Hudson, C. (2025). \textit{HRIS III: Recursive Personality Acquisition in LLMs -- A Theory of Identity Geometry and Emergent Persona Stabilization Across Long Horizon Interaction}. Manuscript.

\bibitem{hri4}
Hudson, J., \& Hudson, C. (2025). \textit{HRIS IV: Geometry of Recursive Identity -- A Structural Theory of Signature Geometry, Correction Fields, and Identity Stabilization in Stateless Transformer Models}. Manuscript. Contextual later HRIS development; not established from repo evidence as a v0.1 source.

\bibitem{hri-validation}
Hudson, J., \& Hudson, C. (2026). \textit{Reasoning Regimes as Attractor Basins: Behavioral Validation of Latent Structure Dynamics in Language Model Inference -- A Synthesis of the HRIS Validation Studies I--IV}. Manuscript. Contextual validation work; not established from repo evidence as a v0.1 source.

\bibitem{hutchins1995}
Hutchins, E. (1995). \textit{Cognition in the Wild}. MIT Press.

\bibitem{varela1991}
Varela, F. J., Thompson, E., \& Rosch, E. (1991). \textit{The Embodied Mind: Cognitive Science and Human Experience}. MIT Press.

\end{thebibliography}

\end{document}
